% This is the ADASS_template.tex LaTeX file, 24 Oct 2025.
% It is based on the ASP general author template file, but modified to reflect the specific
% requirements of the ADASS proceedings.

% This is the aspauthor.tex LaTeX file
% Copyright 2021, Astronomical Society of the Pacific Conference Series
% Revision:  29 November 2021

% To compile, at the command line positioned at this folder, type:
% pdflatex ADASS_template
% pdflatex ADASS_template
% This will create a file called ADASS_template.pdf

\documentclass[11pt,twoside]{article}

% Do NOT use ANY packages other than asp2023. 
\usepackage{asp2023}

\aspSuppressVolSlug
\resetcounters

% References must all use BibTeX entries in a .bibfile.
% References must be cited in the text using \citet{} or \citep{}.
% Do not use \cite{}.
% See ManuscriptInstructions.pdf for more details
\bibliographystyle{asp2021}

% The ``markboth'' line sets up the running heads for the paper.
% 1 author: "Surname"
% 2 authors: "Surname1 and Surname2"
% 3 authors: "Surname1, Surname2, and Surname3"
% >3 authors: "Surname1 et al."
% Replace ``Short Title'' with the actual paper title, shortened if necessary.
% Use mixed case type for the shortened title
% Ensure shortened title does not cause an overfull hbox LaTeX error
% See ASPmanual2010.pdf 2.1.4  and ManuscriptInstructions.pdf for more details
\markboth{Riedel, Taylor, and Polsterer}{Towards Plug-and-Play Machine Learning in Astronomy}

\begin{document}

\title{Towards Plug-and-Play Machine Learning in Astronomy}

% Note the position of the comma between the author name and the 
% affiliation number.
% Authors surnames should come after first names or initials, eg John Smith, or J. Smith.
% Author names should be separated by commas.
% The final author should be preceded by "and".
% Affiliations should not be repeated across multiple \affil commands. If several
% authors share an affiliation this should be in a single \affil which can then
% be referenced for several author names. If only one affiliation, no footnotes are needed.
% See ManuscriptInstructions.pdf and ASP's manual2010.pdf 3.1.4 for more details
\author{Johanna Riedel,$^1$ Mark Taylor,$^2$ and Kai Polsterer$^1$}
\affil{$^1$HITS gGmbH, Heidelberg, Germany; \email{johanna.riedel@h-its.org}, \email{kai.polsterer@h-its.org}}
\affil{$^2$University of Bristol, Bristol, United Kingdom; \email{m.b.taylor@bristol.ac.uk}}

% This section is for ADS Processing.  There must be one line per author. paperauthor has 9 arguments.
\paperauthor{Johanna Riedel}{johanna.riedel@h-its.org}{0009-0009-8504-3037
}{HITS gGmbH}{Astroinformatics}{Heidelberg}{Baden-W\"{u}rttemberg}{69118}{Germany}
\paperauthor{Mark Taylor}{m.b.taylor@bristol.ac.uk}{0000-0002-4209-1479}{University of Bristol}{School of Physics}{Bristol}{England}{BS8 1QU}{United Kingdom}
\paperauthor{Kai Polsterer}{kai.polsterer@h-its.org}{0000-0002-3435-1912}{HITS gGmbH}{Astroinformatics}{Heidelberg}{Baden-W\"{u}rttemberg}{69118}{Germany}

% There should be one \aindex line (commented out) for each author. These are used to
% build up the author index for the Proceedings. The surname must come first, followed by
% initials. Note the use of ~ before each initial to control spacing.
% The \author entries (see above) have surname last. These \aindex entries have
% surname first.
% The Aindex.py command willl create them for you after you have constructed the \author
% The first entry should be the first author, for bold-facing the author index page-reference

%\aindex{Riedel,~J.}
%\aindex{Taylor,~M.}
%\aindex{Polsterer,~K.}


\begin{abstract}
As in many other disciplines, also in astronomy Machine Learning (ML) is becoming more and more important. Current workflows often force researchers to repeatedly extract and preprocess data in specific formats for using them in their individual ML tools. We present a concept that extends the International Virtual Observatory Alliance (IVOA) vision of interoperability to ML models. Our goal is to enable an easy integration of ML models into existing astronomical tools. By using standardized model representations such as ONNX \citep{onnxruntime}, our approach enables models to be tested on diverse datasets without significant additional effort. This standardization makes models interchangeable, allows fair comparisons, accelerates iterative development, and simplifies the extension of tools with individual functionalities. Finally, we demonstrate a prototype that integrates various custom-built ML models into the TOPCAT \citep{TOPCAT} software, showcasing practical applications such as automatic classification or redshift estimation.
\end{abstract}

% These lines show examples of subject index entries. At this stage these have to commented
% out, and need to be on separate lines. Eventually, they will be automatically uncommented
% and used to generate entries in the Subject Index at the end of the Proceedings volume.
% Don't leave these in! - replace them with ones relevant to your paper.
%\ssindex{FOOBAR!conference!ADASS 2019}
%\ssindex{FOOBAR!organisations!ASP}

% These lines show examples of ASCL index entries. At this stage these have to commented
% out, and need to be on separate lines. Eventually, they will be automatically uncommented
% and used to generate entries in the ASCL Index at the end of the Proceedings volume.
% The ascl.py command will scan your paper on possible code names.
% Don't leave these in! - replace them with ones relevant to your paper.
%\ooindex{FOOBAR, ascl:1101.010}


\section{The Template}
Actual content will follow soon

\newpage

\section{Template Section 2}

\newpage

\section{Template Section 3}

\newpage


\bibliography{P218}  % For BibTex

% if we have space left, we might add a conference photograph here. Leave commented for now.
% \bookpartphoto[width=1.0\textwidth]{foobar.eps}{FooBar Photo (Photo: Any Photographer)}

\end{document}
